% 原书表2，原书第320页；逐项对照原书及MinerU Markdown恢复。
\begin{table}[htbp]
\centering\small
\setlength{\parskip}{0pt}
\caption{入藏官兵石刻}\label{tab:tibet-inscriptions}
\renewcommand{\arraystretch}{1.5}
\begin{tabular}{|c|c|c|c|}
\hline
 & 碑题 & 落款时间 & 大约字数 \\
\hline
1 & 功垂百代 & 康熙五十九年十一月 & 680 \\
\hline
2 & 安藏碑记 & 康熙五十九年秋 & 1240 \\
\hline
3 & 用昭万世 & 康熙五十九年十一月 & 770 \\
\hline
4 & 异域流芳 & 康熙六十年春 & 被覆盖 \\
\hline
5 & 流芳百世 & 康熙六十年春 & 被覆盖 \\
\hline
6 & 永载千秋 & 康熙六十一年春 & 被覆盖 \\
\hline
\end{tabular}
\par\nointerlineskip\vspace{3pt}
\begin{minipage}{\linewidth}\small
（资料来源：邓锐龄，《拉萨布达拉山东侧康熙时汉文摩崖考释》，339。）
\end{minipage}
\end{table}
